\documentclass[11pt]{article}
\usepackage[T1]{fontenc}
\usepackage{lmodern}
\usepackage[a4paper,margin=28mm]{geometry}
\usepackage{amsmath,amssymb,amsthm,mathtools,booktabs,mathrsfs}
\usepackage[hidelinks]{hyperref}
\hypersetup{pdftitle={Signed impulse estimates and global smooth solutions for finite multispecies relativistic Vlasov--Maxwell},pdfauthor={Alec Kriebel}}
\newcommand{\R}{\mathbb R}
\newcommand{\Id}{\mathrm{Id}}
\newcommand{\supp}{\operatorname{supp}}
\newcommand{\diver}{\operatorname{div}}
\newcommand{\curl}{\operatorname{curl}}
\newtheorem{theorem}{Theorem}[section]
\newtheorem{proposition}[theorem]{Proposition}
\newtheorem{lemma}[theorem]{Lemma}
\theoremstyle{remark}\newtheorem{remark}[theorem]{Remark}
\numberwithin{equation}{section}
\setlength{\emergencystretch}{2em}
\title{Signed impulse estimates and global smooth solutions for finite multispecies relativistic Vlasov--Maxwell}
\author{Alec Kriebel\\\small ORCID: \href{https://orcid.org/0009-0001-9320-500X}{0009-0001-9320-500X}}
\date{October 6, 2026}
\begin{document}
\maketitle
\begin{abstract}
We extend the signed-impulse proof of the one-species
three-dimensional relativistic Vlasov--Maxwell theorem in OpenAI's family 362
manuscript to a fixed finite collection of particle species. Mass normalization
gives a common velocity law and constant species-dependent accelerations.
The entire source--receiver cancellation has coefficient $e_a e_b/m_a$;
the source acceleration cancels before its sign is discarded. A positive
mass-weighted energy controls a disjoint union of species label spaces,
and the angular occupation and selected receiver coefficient use one
simultaneous momentum bootstrap. We prove the resulting global
classical theorem for nonnegative smooth compact phase data and compatible
$C_b^\infty\cap L^2$ Maxwell fields, without smallness, symmetry or neutrality.
The universal geometric identity is inherited; the new estimates keep all
source and receiver coefficients explicit and close the shared-field bootstrap.
\end{abstract}

\section{Problem, attribution and exact scope}
We use rationalized units with light speed and vacuum permittivity one.
For fixed $N\ge1$, masses $m_a>0$ and charges $e_a\in\R$, physical number
densities $F_a(t,x,p)$ have velocities
$p/(m_a^2+|p|^2)^{1/2}$ and momentum accelerations
$e_a(E+u_a\times B)$. Put
\begin{equation}\label{eq:normalization}
 \begin{gathered}
 v=p/m_a,\qquad f_a(t,x,v)=m_a^3F_a(t,x,m_av),\\
 q(v)=\sqrt{1+|v|^2},\qquad u(v)=v/q(v),\qquad\lambda_a=e_a/m_a.
 \end{gathered}
\end{equation}
Thus $f_a\,dv=F_a\,dp$; $f_a$ remains a number density. The equations are
\begin{equation}\label{eq:system}
\begin{aligned}
 \partial_t f_a+u\cdot\nabla_x f_a+
      \lambda_a(E+u\times B)\cdot\nabla_v f_a&=0,\quad 1\le a\le N,\\
 \partial_t E-\curl B=-j,\qquad \partial_t B+\curl E&=0,\\
 \diver E=\rho,\qquad\diver B&=0,\\
 \rho=\sum_a e_a\int f_a\,dv,\qquad j&=\sum_a e_a\int uf_a\,dv.
\end{aligned}
\end{equation}
There is no background charge. A classical solution has $C^1$ densities and
fields, solves these equations pointwise, has fields continuous into $L^2_x$,
and has compact particle phase support on each compact time interval.

\begin{theorem}[Finite-species global theorem]\label{thm:main}
For every fixed finite species list as above, initial data
\begin{equation}\label{eq:data}
 0\le f_{a0}\in C_c^\infty(\R^6),\qquad
 E_0,B_0\in C_b^\infty(\R^3;\R^3)\cap L^2,
 \quad\diver E_0=\sum_a e_a\int f_{a0}\,dv,\quad\diver B_0=0
\end{equation}
give a unique global classical solution of \eqref{eq:system}. It is smooth on
each finite time interval, its particle phase support is compact there,
and both constraints persist. On a finite maximal lifespan, a bound on
$\max_a\sup_{(x,v)\in\supp f_a(t)}|v|$ is a continuation criterion.
\end{theorem}

$C_b^\infty$ means every spatial derivative is bounded. The field assumption
does not require all derivatives to be square integrable. It permits a nonzero
net charge and its Coulomb tail. Finite field energy supplies cone and
spatial budgets, while bounded initial derivatives control homogeneous terms.
Requiring all derivatives bounded is a convenient smooth-persistence class;
no sharp minimal-regularity assertion is made. Constants below may depend on the data,
the fixed horizon and all fixed masses and charges; no massless or infinite
species limit is asserted. The first new special case is $N=2$,
$m_1=m_2=1$, $e_1=1,e_2=-1$. Noncompact momentum tails, collisions and curved
spacetime are outside the theorem.

The underlying one-species unrestricted large-data theorem and the angular
selection mechanism are due to the OpenAI manuscript \cite{OAI362}, supplied
as part of its October 2026 public repository release, with manuscript date
September 23, 2026. We use its exact pinned version
\texttt{adc7f1241b42e322a6451854ab7e4b4c146bf78a}; the date printed on a
manuscript is not evidence of an earlier public disclosure. This note's
contribution is the signed cross-species estimate and its coupled closure,
not an independent solution of the one-species problem. Glassey--Schaeffer
\cite{GlasseySchaeffer1988} already use mass-normalized momentum and signed
weighted source sums. Their large-data nearly-neutral theorem requires small
signed phase-space density in $C^1$ and small compactly supported $C^2$
fields, rather than arbitrary signed data. Momentum normalization,
multispecies formulation and bounded-momentum continuation are classical;
see Glassey \cite{Glassey} and Luk--Strain \cite{LukStrain}.

If all charged species share a nonzero ratio $\lambda_a=\lambda$, set
$g=\lambda^2\sum_{e_a\ne0}m_af_a$ and
$(\widehat E,\widehat B)=\lambda(E,B)$. Then $g\ge0$ solves the
unit-coefficient one-species system with these fields; its sources are the
transformed Maxwell sources. Charged species can then be transported in
that global field, and neutral species are free. This special case is already
an immediate consequence of the inherited one-species theorem and known
normalization. The extension requiring the estimates below has distinct
charge-to-mass ratios, in particular the two unit-mass charges $+1,-1$.

Recent multispecies work of Mattingly--Pankavich--Ben-Artzi
\cite{MPB2026} sharpens decay and scattering for small globally neutral data.
Its Theorem 1.1 recalls the earlier Glassey--Strauss small-data existence
result; Theorems 1.5--1.6 concern scattering and decay. Those hypotheses do
not duplicate the arbitrary-data, nonneutral theorem considered here.

\section{Positive energy and separate label spaces}
For species $a$, characteristics satisfy
\begin{equation}\label{eq:flow}
 Y_a'=u(V_a),\qquad V_a'=\lambda_a(E+u(V_a)\times B).
\end{equation}
The phase vector field has zero divergence, since $u=\nabla_v q$ and
$\diver_v(u\times B)=0$. Hence every flow preserves volume and transports
$f_{a0}$, preserving positivity, particle number and its $L^\infty$ norm.
Integrating transport gives $\rho_t+\diver j=0$ and propagates both Gauss
constraints. Speeds below one give $\supp_x f_a(t)\subset B(0,R_0+t)$ for a
common initial radius $R_0$.

Define
\begin{equation}\label{eq:energy}
 \mathcal U=\tfrac12(|E|^2+|B|^2)+\sum_a m_a\int qf_a\,dv,
 \qquad\mathcal S=E\times B+\sum_a m_a\int quf_a\,dv.
\end{equation}
Multiplying species transport by $m_aq$ gives work
$m_a\lambda_a E\cdot\int uf_a\,dv=e_a E\cdot\int uf_a\,dv$.
Its sum cancels the Maxwell work, including every charge sign. Thus
$\partial_t\mathcal U+\diver\mathcal S=0$, $|\mathcal S|\le\mathcal U$,
and total energy $\mathcal H_0=\int\mathcal U(0,x)\,dx$ is conserved.
Spatial cutoffs with derivative $O(R^{-1})$ justify this identity because
fields are $C_tL^2$ and particle support is compact on the prefix considered.

On a backward cone $y=x-(t-s)n$, put $r=t-s$ and
$d=1-n\cdot u(v)>0$. Its exact flux density is
\[
 \mathcal U-n\cdot\mathcal S=
 \tfrac12\bigl(|E+n\times B|^2+|n\cdot B|^2\bigr)
       +\sum_b m_b\int qdf_b\,dv.
\]
Integrating the shrinking-ball conservation law bounds its cone integral
by $\mathcal H_0$. In particular, for $G=|E+n\times B|+|n\cdot B|$,
\begin{equation}\label{eq:budgets}
 \int r^2G^2\,dn\,ds\le4\mathcal H_0,\quad
 \int r^2\int qdf_b\,dv\,dn\,ds\le\mathcal H_0/m_b,
 \quad\int qf_b\,dx\,dv\le\mathcal H_0/m_b.
\end{equation}
All are positive bounds. No signed density is substituted for energy.

The positive aggregate $f_*:=\sum_b f_b$ has
$\|f_*\|_\infty\le\sum_b\|f_{b0}\|_\infty$ and
$\int qf_*\le\mathcal H_0\sum_bm_b^{-1}$, with the analogous cone bound.
It generally solves no scalar Vlasov equation. Its label space is the
\emph{disjoint union} of the species label spaces, each with its own flow.
Every labelwise calculation below uses the appropriate component before
summing. This distinction is essential for opposite charges.

\section{The signed source--receiver identity}
Fix a supported receiver of species $a$, denoted $X(t),V_X(t)$, with
$a(t)=u(V_X(t))$, $q_X=q(V_X)$ and $K_a=V_X'$. Use $Y_b(s,z),V_b(s,z)$
for a source label of species $b$. On its retarded branch,
\begin{equation}\label{eq:geometry}
 t=s+r,\quad rn=X(t)-Y_b(s,z),\quad
 d=1-n\cdot u_b,\quad D=1-n\cdot a,\quad \eta=1-a\cdot u_b,
 \quad |n|=1.
\end{equation}
The symbol $\eta$ is geometric, not a charge. Both $d,D$ are positive.
The retarded inverse flow has Jacobian $dz=d\,dy\,dv$: its derivative is
the volume-preserving fixed-time inverse flow times
$\Id-(u_b,V_b')\otimes(n,0)$, whose determinant is $d$.
Differentiating the retarded relation gives
\begin{equation}\label{eq:measure}
 \begin{gathered}
 \frac{dt}{ds}=\frac dD,\qquad
 ds\,f_{b0}(z)dz=D\,dt\,f_b(s,y,v)dy\,dv,\\
 r'=d/D-1,\qquad rn'=N_0=(d/D)(a-n)+n-u_b.
 \end{gathered}
\end{equation}
Strict timelikeness makes each noncolliding branch unique. The set of labels
colliding with a fixed receiver at some equal time is null: a time mesh
of size $\zeta$ covers such labels by $O(\zeta^{-1})$ spatial balls of
volume $O(\zeta^3)$, and a bounded momentum set has finite volume.
Volume preservation and bounded initial density give mass $O(\zeta^2)$.
This argument applies separately to every source species.

Linearity of the wave potentials in Maxwell sources yields the source-$b$
kernel, with $b_b=du_b/ds$,
\begin{equation}\label{eq:kernel}
 \mathscr K_b=\frac{n-u_b}{r^2q_b^2d^3}
       +\frac{(n-u_b)(n\cdot b_b)-db_b}{rd^3},\qquad
 b_b=\frac{\lambda_b}{q_b}(\Id-u_b\otimes u_b)(E+u_b\times B).
\end{equation}
Its magnetic kernel is $n\times\mathscr K_b$ and its source multiplier
is $e_b/(4\pi)$. The receiver force multiplier is consequently
\begin{equation}\label{eq:pair-factor}
 c_{ab}=\lambda_a e_b=\frac{e_ae_b}{m_a},
\end{equation}
not $\lambda_a\lambda_b$. The latter normalization would omit $m_b$.
Initial-cone and homogeneous-data terms are bounded by a datum/horizon
constant, multiplied by the fixed coefficients. This follows from bounded
initial momenta and Kirchhoff's estimate using bounded initial derivatives.

Set $Q=D\Id+n\otimes a$, $H=Q/D$, $g=(u_b-n)/d$, and
$k_0=u_b-\eta n/D$. Then $Hg=k_0/d$. The complete branch-measure
interior pair force, after removing $c_{ab}/(4\pi)$, is
\[
 \mathcal K=-\frac Hr\partial_s^{(u_b)}g-\frac{Hg}{r^2q_b^2d}.
\]
The second term is the transport term and must be retained.

\begin{lemma}[Universal pair cancellation]\label{lem:cancellation}
For arbitrary differentiable timelike source and receiver curves with
the velocity law \eqref{eq:normalization},
\begin{equation}\label{eq:signed}
 \mathcal K=-\left(\frac{k_0}{rd}\right)'
  +\frac{\eta}{r^2D^2}\left(\frac{n}{q_X^2D}-a\right)
  +\frac{n(a_t\cdot k_0)}{rD^2}.
\end{equation}
The identity requires no assumption on the signs or sizes of accelerations.
It therefore holds after multiplication by every signed $c_{ab}$.
\end{lemma}
\begin{proof}
The product rule separates differentiation through $u_b,n,r,a$.
For any variation $A$ of $a$,
$\partial_a k_0[A]=n(A\cdot k_0)/D$; since $a'=a_t d/D$, this gives
the last term. The remaining geometric identities are
\[
 N_0\cdot u_b-dr'-q_b^{-2}=-\eta d/D,
 \qquad rd(k_0)'_n=-\frac{\eta d}{D}
       \left(N_0+\frac{n(N_0\cdot a)}D\right),
\]
\[
 N_0+\frac{n(N_0\cdot a)}D+k_0=
       \frac dD\left(a-\frac{n}{q_X^2D}\right).
\]
They follow by expanding \eqref{eq:geometry}, using
$|u_b|^2=1-q_b^{-2}$ and $|a|^2=1-q_X^{-2}$.
Substitution into the differentiated primitive gives \eqref{eq:signed}.
These identities are also checked by the denominator-cleared universal
polynomial certificate distributed with this note. The proof, rather than
the supplementary finite rational samples, establishes all quantifiers.
\end{proof}

After conversion to $dt\,dy\,dv$, the nondifferentiated residuals are
\begin{equation}\label{eq:residuals}
 \mathcal R_{ab}=c_{ab}\frac{f_b\eta}{r^2D}
       \left(\frac{n}{q_X^2D}-a\right),\qquad
 \mathcal A_{ab}=c_{ab}\frac{f_b n(a_t\cdot k_0)}{rD},\qquad
 a_t=\frac{\Id-a\otimes a}{q_X}K_a.
\end{equation}
Here $K_a$ is the \emph{full} acceleration from the common field.
No additional factor $\lambda_a$ belongs in this last formula.

\section{Uniform occupation and coefficient transfer}
We spell out why the analytic estimates of \cite[Sections 4--8]{OAI362}
apply to this collection, rather than invoking its one-species theorem
on separately forced populations.

Fix a horizon $T_0$ and a compact classical prefix $[0,t_*]$ with
$t_*\le T_0$. Choose a sufficiently large dyadic $P$ such that
\[
 P\ge1024\max_{b,z\in\supp f_{b0},\,t\le t_*}q(V_b(t,z)).
\]
Empty label components are omitted. Write
\[
 L=\log(2+P),\qquad S(w)=P^2L/w.
\]
The simultaneous bootstrap assumes, for every species and supported label,
\begin{equation}\label{eq:bootstrap}
 |V_a(t_2)-V_a(t_1)|\le MP\sqrt I+AI S(w),\qquad I=t_2-t_1,
 \qquad w/8\le q(V_a(t))\le8w\quad(t\in[t_1,t_2]).
\end{equation}
One chooses $M,A$ only after establishing constants independent of them.

Use the exact dyadic partitions of \cite[Section 4]{OAI362}:
$q_b\asymp p$, $q_X\asymp w$, $d\asymp\theta^2$,
$D\asymp\phi^2$, $\eta\asymp\sigma^2$, $r\asymp h$ after normalizing
the radius by $\max(1,T_0)$. Source representatives are $p=64\cdot2^k$;
the fixed comparisons imply $64\le p\le P$, $\theta\ge1/p$,
$\phi\gtrsim1/w$, $\sigma\gtrsim1/m$, $m=\min(p,w)$.
For $\theta<\kappa\phi$ the source/receiver separation scale is
$\sigma\asymp\phi$ and its partition is omitted; otherwise retain
$\sigma$ and put $\nu=\min(\theta,\sigma)$,
$k_1=1+\phi/\sigma$, $\omega=\sigma$.
In the small sector put $\nu=\theta$, $k_1=1$, $\omega=\phi$.
The fixed $\kappa$ is small enough for all enlarged overlapping supports.

For a receiver interval of length at most $I\le(w/P)^2$, define the aggregate
hit measure $\mathsf M=\sum_b\int r^2D N_b\,dt\,ds\,dn$, where
$N_b$ is the bin-restricted velocity integral of $f_b$.
Each species has $N_b\le C_b p^3\nu^2$, by its conserved $L^\infty$ norm
and the velocity cap. At fixed source time the spatial map
$(t,n)\mapsto X(t)-(t-s)n$ is injective with Jacobian $r^2D$.
By \eqref{eq:measure},
\begin{equation}\label{eq:hit}
 \mathsf M=\sum_b\int f_{b0}(z)dz\int_{\text{bin hits}}ds,
 \qquad \mathsf M\le C\max(I,h)/p,
 \qquad\mathsf M\le CI\phi^2/\theta^2.
\end{equation}
The energy bounds in \eqref{eq:budgets} prove the first estimate; branch
time change and the finite sum of species particle numbers prove the second.

For a sufficiently large fixed $C_1$, set
\begin{equation}\label{eq:lambda}
 \ell^{-1/2}=C_1\frac Pm
       \left(\frac M\omega+\sqrt{\frac{AL}{\omega}}\right).
\end{equation}
The first-exit argument using \eqref{eq:bootstrap} bounds source and receiver
increments on cells of length $O(\ell)$ by $\varepsilon m\omega$.
The velocity derivative has norm at most $1/q$, giving velocity changes
$\varepsilon\omega$. This step uses each species' own signed bootstrap;
there is no equality of its acceleration with another species' acceleration.
On actual bins $\omega>C_0/m$, projected relative velocity is bounded
below by $c\omega$ on a cell. The same-time separation at a hit is
$O(h(\phi+\omega))$, so monotone crossing gives duration $Chk_1$ per
label. The label mass at a common cell time is at most $C/p$ after
summing the species energy bounds. Consequently
\begin{equation}\label{eq:cell}
 \mathsf M\le C(hk_1/p)\max(1,I/\ell,h/\ell).
\end{equation}

For the small sector, average the cell mass using the two indicators
$q\asymp p,|u-a(s)|\asymp\phi$ and
$q\asymp p,\sqrt{(1-n_{\rm ret}\cdot a(t_{\rm ret}))/2}\asymp\phi$
over one enlarged source-time interval of length $O(I)$.
Define $\mathcal E_{p\phi}$ with density $qf_*$ and normalization $1/I$.
Both indicator families have bounded overlap and unique retarded receiver
time, so $\sum_{p,\phi}\mathcal E_{p\phi}\le C$.
Full stable cells can be averaged; the at most two clipped cells use their
direct mass bound. When $\ell<h\le I$, the spatial hit map and the second
indicator give $\mathsf M\le CI\mathcal E_{p\phi}/p$ directly.
These arguments yield
\begin{equation}\label{eq:refined}
 \mathsf M\le C(hk_1/p)R^2,\quad
 R=1+\sqrt{I/\ell}\sqrt{\mathcal E},\quad
 \mathcal E=\mathcal E_{p\phi}\ \text{in the small sector},\quad
 \mathcal E=1\ \text{otherwise},
\end{equation}
except the deficit bins $\omega\le C_0/m$ with $h\le\min(I,\ell)$.
On that exception $\ell\le CP^{-2}$ and the direct near bounds suffice.
Every cell and indicator proof operates on each component of the disjoint
label space, then adds positive masses. Thus the scalar transport equation
for $f_*$ is never assumed.

\subsection{Direct bounds, derivatives and uniform constants}
The source acceleration satisfies
\[
 |b_b|\le C|\lambda_b|q_b^{-1}(G+\theta^2|B|),\quad
 |n\cdot b_b|\le C|\lambda_b|\theta q_b^{-1}(G+\theta^2|B|),\quad
 |q_b'|\le C|\lambda_b|(G+\theta|B|).
\]
For every pair the transport majorants are
$C|c_{ab}|f_b/(r^2p^2\theta^2)$ and
$C|c_{ab}|f_b\phi/(r^2p^2\theta^3)$; the good-field and magnetic
majorants are
\begin{equation}\label{eq:maj}
 C|c_{ab}\lambda_b|\frac{f_b\phi G}{rp\theta^2},\qquad
 C|c_{ab}\lambda_b|\frac{f_b\phi|B|}{rp}.
\end{equation}
Perpendicular receiver projection supplies an additional factor $\phi$
to the acceleration-kernel majorants, because the kernel numerator is
tangent to $n$. Taking the finite maximum of
$1,|c_{ab}|,|c_{ab}\lambda_b|$ bounds all these by the same positive
majorants with density $f_*$, up to a fixed constant.

Cone $G^2$ budgets and spatial $B^2$ budgets are shared over the angular
bins, as in \cite[Section 5]{OAI362}. Cauchy--Schwarz, the density cap and
\eqref{eq:refined} give the following representative near/far bounds
in the small sector, suppressing their fixed constants and the square
roots of bounded field budgets:
\[
 T_\perp:\quad \min(Ihp\phi^3/\theta,
                  \mathsf M/(h^2\phi p^2\theta^3)),\qquad
 S_g:\quad\min(I\sqrt h\,p^2\phi^2,
                \sqrt I\sqrt{p\mathsf M}/(\theta h)),
\]
\[
 S_b:\quad\min(\sqrt{I(I+h)}\sqrt h\,p^2\theta^2\phi,
        \sqrt{I+h}\sqrt{p\mathsf M}\theta/(\phi h)).
\]
The remaining sector replaces the relevant $\theta^2$ density factor
by $\nu^2$, and retains $k_1$. These are exactly the nonnegative inputs
to the radial balances and angular sums of \cite[Section 5]{OAI362}.
For example $\sum_h\min(Ah^\alpha,Bh^{-\beta})
\le C_{\alpha,\beta}A^{\beta/(\alpha+\beta)}B^{\alpha/(\alpha+\beta)}$.
The finite species constants multiply $A,B$ outside these balances.
They do not alter any exponent, shared-budget overlap or bootstrap parameter.
The resulting baseline and absolute bounds are
\begin{equation}\label{eq:direct}
 \begin{gathered}
 \text{baseline cost}\le C\mathcal B(I,w),\qquad
 \int_J|K_a|\,dt\le C\sqrt w\,\mathcal B(I,w),\\
 \mathcal B(I,w)=P\sqrt I+(M+\sqrt A)IS(w).
 \end{gathered}
\end{equation}
Here baseline means transport, magnetic remainder and the remaining-sector
good-field cost, as precisely enumerated in that source proposition.
The exceptional deficit near costs have ratio $O(L^3/P)$ to $IS(w)$.

The signed primitive creates source derivatives with coefficient
$|c_{ab}\lambda_b|$ and exactly the majorants \eqref{eq:maj}; geometric
cutoff derivatives have majorant $C|c_{ab}|f_b\phi/(r^2\theta)$.
The central residual is bounded by $C|c_{ab}|f_b/r^2$.
The receiver residual and receiver-cutoff derivatives have majorant
\begin{equation}\label{eq:receiver-majorant}
 C|c_{ab}|\frac{f_b\theta|K_a|}{wr\phi^2}.
\end{equation}
Projection adds $\phi$ to these bounds. For the moving projector
$\Pi=w(\Id-\xi\otimes\xi)/|V_X|$, $\xi=V_X/|V_X|$,
its derivative has the projected receiver size.
A Lipschitz weight derivative contributes
$C|c_{ab}|\,|\psi_t|f_b\theta/(r\phi)$.

Choose representative selections fixed for the receiver interval and common
to every source component. Derivatives of their selected partition sum equal
the negative sum of unselected derivatives, with common physical multiplier
factored out. Supports meeting have representatives within factor four.
Thus source and geometric derivative costs can be assigned to unselected
neighbors as in \cite[Section 6]{OAI362}; no different kernels or charges
are canceled against each other. Integration by parts is valid componentwise.
Initial-cone terms have total cost $C\int|\psi|$, since $Dk_0=Du_b-\eta n$
is bounded and each initial $d$ is bounded below. A radial tip cutoff has
error $O_{\rm prefix}(\gamma)$ and can be removed on the fixed prefix.
Such justification constants do not enter the uniform estimates.

\subsection{Direction occupation and selected coefficients}
Apply the projected identity with selection
$\theta<\kappa\phi$, $p\theta\sqrt h\le1$ to the successive first exits
of direction of a receiver of \emph{any} species. For
$w\ge P^{3/5}$, $w^{-3/5}\le\delta\le w^{-1/10}$ and a time interval
of length at most $(w/P)^2$, the direction count proof of
\cite[Section 7]{OAI362} uses only \eqref{eq:bootstrap},
\eqref{eq:refined}, \eqref{eq:direct} and the projected majorants above.
The total receiver/projector coefficient is $CL^2/w$; the time-weight
cost is $CL^2$ per exit. Bulk terms are integrated over disjoint intervals,
so they share one length bound rather than gaining an exit count.
The estimate is
\[
 \tfrac12N_cw\delta\le C(M,A)wL^6+CL^2N_c.
\]
Absorb the last term using $w\delta\ge P^{6/25}$ to obtain
$N_c+1\le C(M,A)L^6/\delta$. This uses no improved occupation input.

Separate source and receiver direction partitions then give, in the window
\begin{equation}\label{eq:window}
 p,w\ge P^{7/10},\qquad
 j^{-49/100}\le\phi\le j^{-4/25}\quad(j=p,w),
\end{equation}
\begin{equation}\label{eq:improved}
 \mathsf M\le C(h/p)\phi^{-29/25}
     \max(1,I/I_m^0,h/I_m^0),\qquad I_m^0=(m/P)^2.
\end{equation}
Indeed, with stable energy cells of length
$I_m^0/[C_2(M^2+AL)]$, bad hits crossing a receiver direction boundary
have source duration $O(h)$ per boundary. Good hits have same-time
separation $O(h\phi)$ and monotone relative velocity $\Omega(\phi)$ on
the common refinement of source and receiver direction cells. The refinement
count is their sum, so duration is $C(M,A)hL^6/\phi$ per label.
Summing positive species masses and cells gives the pre-absorption cost
$C(M,A)(h/p)\phi^{-1}L^7\max(1,I/I_m^0,h/I_m^0)$.
The window gives $\phi^{-4/25}\ge P^{56/3125}$, absorbing the fixed
$C(M,A)L^7$ for a sufficiently large threshold. The final $C$ in
\eqref{eq:improved} is independent of $M,A$. This is the required order
of dependency, rather than a circular use of improved occupation to count
directions.

Set $\epsilon=10^{-5}$, $\beta=29/25$ and
$W=(P\theta)^{-\epsilon}\phi^\epsilon h^\epsilon$.
Let $\mathsf M_*$ be the minimum of the \emph{pure} entries
\[
 \frac{\max(I,h)}p,\quad
 \frac hp\max(1,I/\ell,h/\ell)\ (\phi>C_0/m),\quad
 \frac hp\phi^{-\beta}\max(1,I/I_m^0,h/I_m^0)\
       \text{when \eqref{eq:window} holds}.
\]
Then $\mathsf M\le C\mathsf M_*$, with the constant outside the entries.
Select a small-sector representative if
\begin{equation}\label{eq:selection}
 \min\left(I\sqrt h\,p^2\phi^2,
      \frac{\sqrt I\sqrt{p\mathsf M_*}}{\theta h}\right)>P\sqrt I\,W.
\end{equation}
Unselected good-field costs sum to $CP\sqrt I$ because $\sum W\le C$.
At fixed receiver time, \eqref{eq:receiver-majorant}, the density cap and
the positive cone budget give aggregate coefficient at most $CU$, where
\begin{equation}\label{eq:U}
 U=w^{-1}\min\{h^2(p\theta)^3,(p\theta h\phi^2)^{-1}\}.
\end{equation}
Both estimates are available with fixed species-dependent constants. Taking
their maximum bounds the minimum by $CU$. The exact representative
argument of \cite[Section 8]{OAI362} gives
\begin{equation}\label{eq:safe}
 U\le w^{-1/2}W\qquad\text{on selected representatives for }P\ge P_*(M,A).
\end{equation}
To check its decisive compatibility, write
$p=P^{1-\alpha}$, $w=P^z$, $c=1-z$, $p\theta=P^b$,
$\phi=P^{-y}$, $h=P^{-H}$ and
$\Delta=\epsilon(\alpha+b+y+H)$.
Near selection gives $2\alpha+H/2+2y<z+\Delta$.
Energy far selection and unsafety force
$\Delta<4z/99995<.001z$,
$H>.496z+3(\alpha+c)$ and actual separation $m\phi>P^{.3}$.
Stable-cell far selection then forces the window \eqref{eq:window}.
Improved occupation and unsafety imply
\[
 \beta y>H/3+z/3+2\min(\alpha,c)-8\Delta/3,
 \qquad\beta y<29z/50-29H/100+(29/50)\Delta.
\]
Consequently $H/z<37487/93500<62/125$, contradicting $H/z>.496$.
The intervening window inequalities and exact rational margin are checked
by the accompanying selection certificate and reproduced in the audit
supplement; no species coefficient occurs inside these pure inequalities.

Summing \eqref{eq:safe} bounds the coefficient of the \emph{full} $|K_a|$
by $Cw^{-1/2}$, uniformly over all receivers. It cancels the $\sqrt w$
loss in \eqref{eq:direct}. The spatial, central and source derivative
balances and endpoint-weight estimates of \cite[Section 8]{OAI362} use
exactly the transferred majorants listed above. Hence, for an endpoint-zero
weight with ramps of length $\eta_0 I$,
\begin{equation}\label{eq:weighted}
 \left|\int_J\psi K_a\,dt\right|
 \le C_{\eta_0}\bigl[P\sqrt I+(M+\sqrt A)IS(w)\bigr].
\end{equation}
The constant depends on the fixed species data but not on $M,A$.
This proves parameter robustness of the analytic chain: positivity enters
only the aggregate majorants and energy, while every signed kernel identity
was first taken pairwise with its full fixed coefficient.

\section{Coupled closure and exclusion of finite-time blowup}
Let $\mathcal B_w(I)=MP\sqrt I+AI S(w)$.
If it is at least $32w$, the endpoint energy range directly gives
$|\Delta V_a|\le16w\le\mathcal B_w(I)/2$.
Otherwise $M\ge32$ implies $I\le(w/P)^2$.
Restore the two edge ramps in \eqref{eq:weighted} by integrating the
removed weight against $V_a(t)-V_a(t_1)$ on the left and
$V_a(t)-V_a(t_2)$ on the right. Each monotone ramp has derivative norm one,
so the simultaneous signed bootstrap, rather than an absolute force bound,
gives total cost $2\mathcal B_w(\eta_0I)$. Thus
\[
 |\Delta V_a|\le(C_{\eta_0}+2M\sqrt{\eta_0})P\sqrt I+
  [C_{\eta_0}(M+\sqrt A)+2A\eta_0]IS(w).
\]
Choose $2\sqrt{\eta_0},2\eta_0\le1/8$, then
$M\ge\max(32,8C_{\eta_0})$, then
$A\ge\max(1,16C_{\eta_0}M,256C_{\eta_0}^2)$, and finally $P_0$
large enough for all preceding estimates. This gives a strict half-bootstrap
for \emph{every} species with one set of constants.

The set of times up to which all eligible intervals and all species satisfy
\eqref{eq:bootstrap} is closed by truncation and continuity. On a fixed
compact classical prefix, the finite union of supported label sets has a
finite acceleration bound $C_*$. Every interval of sufficiently small length
obeys $C_*I\le MP\sqrt I$. The strict improvement on a longer interval's
prefix and $C_*\varepsilon$ on its suffix prove relative right openness,
uniformly over the finite species list. This temporary $C_*$ affects only
the openness argument; it never enters $M,A,P_0$.
Therefore the signed increment estimate is unconditional on every compact
classical prefix, with constants independent of that prefix's endpoint.

If all particle data vanish, smooth vacuum Maxwell evolution is global.
Otherwise suppose a maximal solution had finite lifespan and unbounded joint momentum.
Let $Q(t)=\max_{a,z\in\supp f_{a0}}q(V_a(t,z))$ and $t_n$ its first hit
of $2^n$ for large $n$. Choose a maximizing species and label at $t_n$,
and $s_n$ its last hit of $2^{n-1}$. Then $s_n\ge t_{n-1}$ even if the
maximizer changes species, and the chosen label remains in
$[2^{n-1},2^n]$ on $[s_n,t_n]$. Apply the signed estimate with
$P=1024\cdot2^n$, $w=2^n$, $\Delta_n=t_n-s_n$. Since $q$ is
one-Lipschitz,
\[
 2^{n-1}\le1024M2^n\sqrt{\Delta_n}+
       1024^2A2^n\Delta_n\log(2+1024\cdot2^n).
\]
One term must exceed $2^{n-2}$, so
$t_n-t_{n-1}\ge\Delta_n\ge c/\log(2+1024\cdot2^n)$.
The sum diverges. Finite-time momentum blowup is impossible.

\section{Local theory, uniqueness and nonneutral continuation}
For Sobolev fields the conventional bounded-momentum continuation input is
Luk--Strain \cite[Theorem 1.1, Footnote 1, Remark 1.2]{LukStrain}:
nonnegative compact phase data in $H^5$, fields in $H^5$, Gauss constraints,
and a uniform finite momentum radius imply continuation. Their remark
explicitly extends the method to finitely many species. Glassey's book
\cite[Chapter 5]{Glassey} formulates the masses and signs and explains
the species-specific streaming operators and separate iterate differences.
Its printed theorem includes neutrality; we do not silently remove that
hypothesis from its statement.

The finite-species Sobolev adaptation uses the same bounded-radius velocity
law after \eqref{eq:normalization}. Source field representations acquire
$e_b\lambda_b$ and receiver force equations acquire $\lambda_a$;
absolute estimates replace them by fixed constants and sum finitely many
species. At bounded radius $R$, all cone divisors obey
$1-n\cdot u\ge1-R/\sqrt{1+R^2}>0$. The division-lemma and differentiated
kernel estimates in \cite{BGP,LukStrain} therefore have the same bounded
constants and logarithmic derivative inequality, with
$\mathcal N(t)=\max_{s\le t}\sum_a\|\nabla_{x,v}f_a(s)\|_\infty$:
\[
 \|(E,B)(t)\|_{W^{1,\infty}}\le C[1+\log(2+\mathcal N(t))],\qquad
 \mathcal N(t)\le\mathcal N(0)+C\int_0^t
       [1+\log(2+\mathcal N(s))]\mathcal N(s)\,ds.
\]
The Osgood integral diverges, controlling derivatives; higher Sobolev norms
then satisfy the usual tame transport--Maxwell Gronwall estimate.
Local $H^6$ iteration uses
$\|j\|_{H^k_x}\le |B_R|^{1/2}\sum_a|e_a|\|f_a\|_{H^k}$ and
the divergence-free phase transport. It gives a common local lifespan,
positivity and smooth persistence for smooth data. The constants depend on
$R$; these continuation estimates do not replace the a priori signed proof.

To reach \eqref{eq:data}, fix any finite horizon $T$ and choose
$L_0>R_0+2T$. The signed smooth compact initial density $\rho_0$ has
Coulomb field $E_C=-\nabla(-\Delta)^{-1}\rho_0$.
It is smooth and belongs to every $H^k$ in three dimensions:
its Fourier transform is $-i\xi\widehat\rho_0/|\xi|^2$, whose
$O(|\xi|^{-1})$ behavior is square integrable near zero, and
$\widehat\rho_0$ decays rapidly at infinity.
This argument permits $\int\rho_0\ne0$.

Both $G_E=E_0-E_C$ and $G_B=B_0$ are smooth divergence-free fields.
For either such $G$, the smooth potential
\[
 A_G(x)=\left(\int_0^1 tG(tx)\,dt\right)\times x
\]
satisfies $\curl A_G=G$, as follows by differentiating $t^2G(tx)$ and
using $\diver G=0$. Let a compact cutoff $\chi$ equal one on $B(0,L_0)$.
Set $\widetilde E_0=E_C+\curl(\chi A_{G_E})$ and
$\widetilde B_0=\curl(\chi A_{G_B})$.
These fields lie in every $H^k$, preserve both constraints and coincide
with the originals on $B(0,L_0)$.

The difference is divergence-free smooth $L^2$ data for vacuum Maxwell.
Its global vacuum evolution $W$ is $C_tL^2$, smooth with bounded derivatives,
and vanishes on $|x|<L_0-t$ by finite propagation. For $t\le T$ it
vanishes throughout every particle region $|x|\le R_0+t$.
Adding $W$ to a solution for modified data, or subtracting it from an
original-data solution, preserves all species transport equations on
their supports, Maxwell equations and constraints, without changing
momenta. Thus bounded-momentum continuation and smooth local theory
transfer to \eqref{eq:data} on any prescribed horizon.

Uniqueness holds in the full classical class, not only the Sobolev construction.
On a common compact interval choose a common momentum radius $R$ and
subtract two solutions. Transport difference energies are bounded by
$C\sum_a\|\delta f_a\|_2\|\delta(E,B)\|_2$, since the other solution's
first density derivatives are bounded on compact phase support. Maxwell's
difference current satisfies
$\|\delta j\|_2\le|B_R|^{1/2}\sum_a|e_a|\|\delta f_a\|_2$.
Spatial cutoffs justify Maxwell energy without assuming $L^2$ field
derivatives. Gronwall applied to
$\sum_a\|\delta f_a\|_2^2+\|\delta(E,B)\|_2^2$ gives equality.
Together with the signed estimate and momentum-doubling argument this
completes the proof of Theorem~\ref{thm:main}. The dependency ledger and
independent adversarial audit records describe the validation performed.

\section{Checks, limitations and disclosure}
If a receiver has $e_a=0$, its momentum is constant and all pair factors
vanish. A source with $e_b=0$ has free transport and no Maxwell contribution;
its positive energy may be retained harmlessly. Identical mass/charge species
can be combined by adding number densities. For a single species with
$m=e=1$ every displayed formula reduces to \cite{OAI362}.
For two unit masses and charges $+1,-1$, same-sign pair coefficients are
$+1$ and cross-sign coefficients $-1$; source accelerations are signed
before exact cancellation. An extreme fixed mass ratio only enlarges
$m_b^{-1}$, $|\lambda_b|$ and $|c_{ab}|$ and the normalized initial norms.
There is no uniform estimate as a mass approaches zero. Neutrality is used
nowhere in the positive energy, kernels, occupation or localization proof.

The assumptions of bounded smooth field derivatives provide the homogeneous
field and vacuum regularity needed by this proof. Arbitrary smooth finite
energy fields without those derivative bounds are not covered. Compact
particle phase support is a hypothesis of the analytic chain. The note
does not infer control of support from energy alone, or a scalar Vlasov
equation for $\sum_b f_b$.

The source package contains exact source hashes, dependency and priority
ledgers, scoped agent audits, and standard-library kernel and exponent
certificates. Those certificates check algebra only. This multispecies
theorem is not claimed to be Lean formalized; a source-library static scan
or existence of comparator files is not formal certification.
AI tools were used extensively in the research, drafting and verification.
Automated adversarial reviews do not constitute conventional human peer
review. This preprint has not undergone conventional human peer review or
refereeing. The ordinary analytic proof is the basis for the theorem; no
reproduced Lean certificate is claimed.

\begin{thebibliography}{9}
\bibitem{OAI362} OpenAI.
\emph{Global classical solutions of the three-dimensional relativistic
Vlasov--Maxwell system}. OpenAI Math Release preprint, manuscript dated
September 23, 2026. Pinned repository input:
\href{https://github.com/openai/math/tree/adc7f1241b42e322a6451854ab7e4b4c146bf78a}{\texttt{adc7f1241b42e322a6451854ab7e4b4c146bf78a}}, family 362.
The manuscript-specific citation is supplied in that family's README.
\bibitem{Glassey} R. T. Glassey. \emph{The Cauchy Problem in Kinetic Theory}.
SIAM, 1996, Chapters 5--6. \href{https://doi.org/10.1137/1.9781611971477}{doi:10.1137/1.9781611971477}.
\bibitem{LukStrain} J. Luk and R. M. Strain.
A new continuation criterion for the relativistic Vlasov--Maxwell system.
\emph{Commun. Math. Phys.} 331 (2014), 1005--1027.
\href{https://arxiv.org/abs/1406.0165v1}{arXiv:1406.0165v1};
\href{https://doi.org/10.1007/s00220-014-2108-8}{doi:10.1007/s00220-014-2108-8}.
\bibitem{GlasseySchaeffer1988} R. T. Glassey and J. Schaeffer.
Global existence for the relativistic Vlasov--Maxwell system with nearly
neutral initial data. \emph{Commun. Math. Phys.} 119 (1988), 353--384.
\href{https://doi.org/10.1007/BF01218078}{doi:10.1007/BF01218078}.
\bibitem{BGP} F. Bouchut, F. Golse and C. Pallard.
Classical solutions and the Glassey--Strauss theorem for the 3D
Vlasov--Maxwell system. \emph{Arch. Ration. Mech. Anal.} 170 (2003), 1--15.
\href{https://arxiv.org/abs/math/0301175}{arXiv:math/0301175}.
\bibitem{MPB2026} G. Mattingly, S. Pankavich and J. Ben-Artzi.
Sharp time decay and scattering of small data solutions to the relativistic
Vlasov--Maxwell system. \href{https://arxiv.org/abs/2609.29013v1}{arXiv:2609.29013v1},
submitted September 24, 2026.
\end{thebibliography}
\end{document}
